\section[Réévaluation]{Réévaluation}
\label{sec:reevaluation-base-regles-reevaluation}

Ici, je présente la méthode de réévaluation.

J'explique que je fais ça à partir du nom persistant d'une orbite (arbre
d'évaluation) et ce pour chacune des orbites identifiées dans la spec à
réévaluer.

\subsection{Reconstruction à partir d'un DAG d'évaluation}
\label{sec:reevaluation-base-regles-reevaluation-reconstruction}

(du début jusqu'à la fin)
Puisque les DAGs de réévaluation sont utilisés pour l'appariement de paramètres
topologiques, ceux-ci ne contiennent non pas des nœuds de règles, mais des brins
issus de l'objet en cours de reconstruction.

\subsection{Stratégies de réévaluation}
\label{sec:reevaluation-base-regles-reevaluation-strategies}

\begin{itemize}
	\item Scission d'un paramètre topologique : évaluation de tout ou partie des paramètres éligibles
	\item Suppression d'un paramètre topologique : interruption de la réévaluation pour ce paramètre (potentiellement de l'objet)
\end{itemize}


%%% Local Variables:
%%% mode: LaTeX
%%% TeX-master: "../../../main"
%%% End:
